\section{Finite Fields From First Principles}

A field is a set in which one can add, subtract, multiply, and divide by any
nonzero element, with the usual algebraic rules.  Familiar examples include the
rational numbers and real numbers.  Cryptographic proof systems usually use
finite fields: fields with only finitely many elements.

The target implementation uses a binary field with $2^{128}$ elements,
\[
  \F_{2^{128}}
  =
  \F_2[X]/(X^{128}+X^7+X^3+X+1).
\]
This means an element is a polynomial of degree at most $127$ with coefficients
in $\{0,1\}$, and multiplication is reduced modulo the irreducible polynomial
$X^{128}+X^7+X^3+X+1$.  Addition is coefficient-wise addition modulo two, which
is bitwise XOR in the usual machine representation.

For the theory, the exact field is less important than three properties:

\begin{enumerate}[label=\textbf{F\arabic*.}]
  \item We can perform field arithmetic exactly.
  \item The field is large enough that random field challenges make cheating
        unlikely.
  \item Linear equations and low-degree polynomial identities can be tested by
        evaluating them at random points.
\end{enumerate}

The protocol relies heavily on linearity.  If $A$ is a linear map, then
\[
  A(\alpha x+\beta y)=\alpha A(x)+\beta A(y).
\]
This simple equation is the reason Blaze can collapse $t$ committed row words
into one openable tensor $c_\star$ and ask the backend to prove one RAA
relation about it.

